% !TEX program = xelatex
% !TEX builder = latexmk
% !Mode:: "TeX:UTF-8:Hard"
%%%%%%%%%%%%%%%%%%%%%%%%%%%%%%%%%%%%%%%%%%%%%%%%%%%%%%%%%%%%%
%% BAC01
%%
%% Andreas Landgraf (2015)
%%
%% see "How to Write a Thesis in LaTeX"
%%     https://www.sharelatex.com/blog/2013/08/02/thesis-series-pt1.html
%%%%%%%%%%%%%%%%%%%%%%%%%%%%%%%%%%%%%%%%%%%%%%%%%%%%%%%%%%%%%

\def\campusSubject{BACHELOR THESIS}
\def\campusTitle{Analysis and Comparison of Wireless Personal Area Networking Technologies for low-power Smart Home Applications}
\def\campusAuthor{Andreas Landgraf}
\def\campusPersKnz{1310317046}
\def\campusCity{Graz}
\def\campusDate{\today}
\def\campusKeywords{WPAN, Smart Home, Bluetooth, Bluetooth Smart, UWB, ZigBee, 6LoWPAN, Wi-Fi, 802.11ad}

\documentclass[
  headexclude,
  footexclude,
  10pt,
  a4paper,
  oneside,
  captions=tableheading,
  bibliography=totoc,
  listof=totoc
]{scrbook}

\RequirePackage[
  automark,
  headsepline
]{scrlayer-scrpage}

\RequirePackage{geometry}
\RequirePackage{layouts}
\RequirePackage{calc}
\RequirePackage{setspace}
\RequirePackage{titlesec}
\RequirePackage[austrian,american]{babel}
\RequirePackage[utf8]{inputenc}
\RequirePackage[T1]{fontenc}
\RequirePackage{fontspec}
\RequirePackage[hyphens]{url}
\RequirePackage{csquotes}
\RequirePackage{amsmath,amsthm,amssymb}
\RequirePackage{color}
\RequirePackage[dvipsnames,svgnames,table]{xcolor}
\RequirePackage{booktabs}
\RequirePackage{textcomp}
\RequirePackage{graphicx}
\RequirePackage{tabularx}
\RequirePackage{multirow}
\RequirePackage[font=small,format=hang,labelfont=bf]{caption}
\RequirePackage{subcaption}
\RequirePackage[colorinlistoftodos,obeyDraft]{todonotes}
\RequirePackage[toc,acronym,style=long,nonumberlist]{glossaries}

\RequirePackage[
  backend=biber,
  style=campus02,
  natbib=true,
  sorting=nty,
  sortlocale=en_US,
  autocite=footnote,
  related=true,
  backref=false,
  hyperref=auto,
  isbn=false,
  url=false,
  citereset=none,
  maxcitenames=3,
  maxbibnames=100,
  block=none,
  citecounter=true,
]{biblatex}

\RequirePackage{enumitem}

\RequirePackage{footnote}
\RequirePackage[multiple]{footmisc}
\RequirePackage{chngcntr}

\RequirePackage[
  pdftitle={\campusTitle},
  pdfauthor={\campusAuthor},
  pdfsubject={\campusSubject, \campusDate},
  pdfkeywords={\campusKeywords},
  pdffitwindow=true,
  pdfpagelayout={OneColumn},
  bookmarks=true,
  bookmarksnumbered=true,
  bookmarksopen=true,
  hyperindex=true,
  colorlinks=false,
  linkcolor=black,
  urlcolor=black,
  citecolor=black,
  unicode=true,
  hidelinks,
  hyperfootnotes=false
]{hyperref}
\RequirePackage{nameref}
\RequirePackage{cleveref}

%%%%%%%%%%%%%%%%%%%%%%%%%%%%%%%%%%%%%%%%%%%%%%%%%%%%%%%%%%%%%
% Bibliographies
%%%%%%%%%%%%%%%%%%%%%%%%%%%%%%%%%%%%%%%%%%%%%%%%%%%%%%%%%%%%%

\counterwithout{footnote}{chapter}

\DeclareLanguageMapping{american}{american-apa}
\NewBibliographyString{retrieved}
\NewBibliographyString{from}
\DefineBibliographyStrings{american}{
  retrieved = {},
  from = {},
}

\addbibresource{literature/mendeley.bib}
\addbibresource{literature/internet.bib}

\defbibfilter{printed}{
  \type{book} \or
  \type{phdthesis} \or
  \type{article} \or
  \type{incollection} \or
  \( \type{techreport} \and \not \keyword{RFC4944} \) \or
  \type{inproceedings}
}

\defbibfilter{online}{
  \type{online} \or
  \( \type{techreport} \and \keyword{RFC4944} \)
}

\newenvironment{bibhyphenation}
{\hyphenpenalty=20
  \tolerance=9999
  \exhyphenpenalty=0
  \linepenalty=0
  \setlength{\emergencystretch}{3em}}
{}

\crefname{table}{table}{tables}
\Crefname{table}{Table}{Tables}
\crefname{figure}{figure}{figures}
\Crefname{figure}{Figure}{Figures}
\crefname{section}{section}{sections}
\Crefname{section}{Section}{Sections}
\crefname{chapter}{chapter}{chapters}
\Crefname{chapter}{Chapter}{Chapters}

%%%%%%%%%%%%%%%%%%%%%%%%%%%%%%%%%%%%%%%%%%%%%%%%%%%%%%%%%%%%%

\geometry{
  left=2.5cm,
  right=2cm,
  top=2.5cm,
  bottom=2.2cm,
}

\graphicspath{{images/}}
\setmainfont{Arial}

\makeglossaries

% Acronyms
\newacronym{wban}{WBAN}{Wireless Body Area Network}
\newacronym{wpan}{WPAN}{Wireless Personal Area Network}
\newacronym{wlan}{WLAN}{Wireless Local Area Network}
\newacronym{wcan}{WCAN}{Wireless Campus Area Network}
\newacronym{wman}{WMAN}{Wireless Metropolitan Area Network}
\newacronym{wwan}{WWAN}{Wireless Wide Area Network}
\newacronym{wran}{WRAN}{Wireless Regional Area Network}
\newacronym{wgan}{WGAN}{Wireless Global Area Network}
\newacronym{han}{HAN}{Home Area Network}
\newacronym{ict}{ICT}{Information and Communications Technology}
\newacronym{iot}{IoT}{Internet of Things}
\newacronym{itu}{ITU}{International Telecommunication Union}
\newacronym{rtd}{RTD}{Research and Technical Development}
\newacronym{sms}{SMS}{Social Media Services}
\newacronym{usb}{USB}{Universal Serial Bus}
\newacronym{wusb}{WUSB}{Certified Wireless USB}
\newacronym{nb}{NB}{Narrowband}
\newacronym{uwb}{UWB}{Ultra Wideband}
\newacronym{hbc}{HBC}{Human Body Communications}
\newacronym{rfid}{RFID}{Radio Frequency Identification}
\newacronym{nfc}{NFC}{Near Field Communication}
\newacronym{pc}{PC}{Personal Computer}
\newacronym{pda}{PDA}{Personal Digital Assistant}
\newacronym{ism}{ISM}{Industrial, Scientific and Medical}
\newacronym{siso}{SISO}{Single-Input Single-Output}
\newacronym{mimo}{MIMO}{Multiple-Input Multiple-Output}
\newacronym{mumimo}{MU-MIMO}{Multi-User MIMO}
\newacronym{dsss}{DSSS}{Direct Sequence Spread Spectrum}
\newacronym{ofdm}{OFDM}{Orthogonal Frequency-Division Multiplexing}
\newacronym{wimax}{WiMAX}{Worldwide Interoperability for Microwave Access}
\newacronym{bwa}{BWA}{Broadband wireless access}
\newacronym{gsm}{GSM}{Global System for Mobile Communications}
\newacronym{tdma}{TDMA}{Time Division Multiple Access}
\newacronym{cdma}{CDMA}{Code Division Multiple Access}
\newacronym{gprs}{GPRS}{General Packet Radio Service}
\newacronym{edge}{EDGE}{Enhanced Data Rates for GSM Evolution}
\newacronym{umts}{UMTS}{Universal Mobile Telecommunications System}
\newacronym{hspa+}{HSPA+}{Evolved High-Speed Packet Access}
\newacronym{wwww}{WWWW}{World Wide Wireless Web}
\newacronym{http}{HTTP}{Hypertext Transfer Protocol}
\newacronym{when}{WHEN}{Wireless Home Entertainment Network}
\newacronym{whan}{WHAN}{Wireless Home Automation Network}
\newacronym{rf}{RF}{Radio Frequency}
\newacronym{m2m}{M2M}{Machine-to-Machine}
\newacronym{h2m}{H2M}{Human-to-Machine}
\newacronym{mih}{MiH}{Machine-in-Human}
\newacronym{whn}{WHN}{Wireless Home Network}
\newacronym{mac}{MAC}{Media Access Control}
\newacronym{l2cap}{L2CAP}{Logical Link Control and Adaption Protocol}
\newacronym{rfcomm}{RFCOMM}{Radio Frequency Communication}
\newacronym{tcs}{TCS}{Telephony Control Protocol Specification}
\newacronym{sdp}{SDP}{Service Discovery Protocol}
\newacronym{hci}{HCI}{Host Control Interface}
\newacronym{fhss}{FHSS}{Frequency Hopping Spread Spectrum}
\newacronym{gfsk}{GFSK}{Gaussian Frequency-Shift Keying}
\newacronym{fsk}{FSK}{Frequency-Shift Keying}
\newacronym{ble}{BLE}{Bluetooth Low Energy}
\newacronym{att}{ATT}{Attribute Protocol}
\newacronym{smp}{SMP}{Security Manager Protocol}
\newacronym{gatt}{GATT}{Generic Attribute Profile}
\newacronym{gap}{GAP}{Generic Access Profile}
\newacronym{pdu}{PDU}{Protocol Data Unit}
\newacronym{crc}{CRC}{Cyclic Redundancy Check}
\newacronym{single_uwb}{DS UWB}{Direct Sequence UWB}
\newacronym{multi_uwb}{MB-OFDM UWB}{Multiband Orthogonal Frequency Division Multiplexing UWB}
\newacronym{ffd}{FFD}{Full-Function Device}
\newacronym{rfd}{RFD}{Reduced-Function Device}
\newacronym{zdo}{ZDO}{ZigBee Device Object}
\newacronym{aps}{APS}{Application Support SubLayer}
\newacronym{6lowpan}{6LoWPAN}{IPv6 over Low-Power Wireless Personal Area Network}
\newacronym{lowpan}{LoWPAN}{Low-Power Wireless Personal Area Network}
\newacronym{aes}{AES}{Advanced Encryption Standard}
\newacronym{pan}{PAN}{Personal Area Network}
\newacronym{ip}{IP}{Internet Protocol}
\newacronym{ipv6}{IPv6}{Internet Protocol version 6}
\newacronym{udp}{UDP}{User Datagram Protocol}
\newacronym{icmp}{ICMP}{Internet Control Message Protocol}
\newacronym{tcp}{TCP}{Transmission Control Protocol}
\newacronym{pbss}{PBSS}{Personal Basic Service Set}
\newacronym{dsl}{DSL}{Digital Subscriber Line}
\newacronym{rfc}{RFC}{Request for Comments}

\renewcommand*{\glspostdescription}{}

\setcounter{secnumdepth}{3}
\setcounter{tocdepth}{3}

\clubpenalty = 10000
\widowpenalty = 10000
\displaywidowpenalty = 10000

\frenchspacing
\renewcommand{\singlespacing}{\setstretch{1.09}}
\renewcommand{\onehalfspacing}{\setstretch{1.45}}
\singlespacing
\newlength{\singlebaselineskip}
\setlength{\singlebaselineskip}{\baselineskip}
\onehalfspacing
\renewcommand{\arraystretch}{1.09}
\addtokomafont{sectioning}{\setstretch{1}\selectfont}

\renewcommand*{\chapterheadstartvskip}{\vspace*{4mm}}
\renewcommand*{\chapterheadendvskip}{\vspace*{0mm}}

\titlespacing*{\section}{0mm}{4mm}{0mm}
\titlespacing*{\subsection}{0mm}{4mm}{0mm}

\setlength{\parskip}{3mm}
\setlength{\itemsep}{1em}

\renewcommand{\chaptermark}[1]{\markboth{#1}{}}

% clear header and footer
\clearscrheadfoot
\ihead[]{\leftmark}
\chead[]{}
\ohead[]{}
\ifoot[]{}
\cfoot[]{}
\ofoot[]{\thepage}

\setkomafont{title}{\fontspec{Arial}\Large\bfseries}
\setkomafont{pageheadfoot}{\fontspec{Arial}\normalsize}
\setkomafont{chapter}{\fontspec{Arial}\Large\bfseries}
\setkomafont{section}{\fontspec{Arial}\Large\bfseries}
\setkomafont{subsection}{\fontspec{Arial}\large\bfseries}
\setkomafont{chapterentry}{\fontspec{Arial}\bfseries}
\addtokomafont{chapter}{\MakeUppercase}

%%%%%%%%%%%%%%%%%%%%%%%%%%%%%%%%%%%%%%%%%%%%%%%%%%%%%%%%%%%%%

\renewcommand*{\maketitle}{
  \clearpage
  \thispagestyle{empty}
  \null
  \begin{center}
    {\textbf{\campusSubject} \par}
    \vskip 2.5em
    {\hrule\vskip 2em \titlefont\huge\campusTitle\vspace*{1em}\hrule\vspace*{1.5em}\par}
    {accomplished at \par}
    \vskip .6em
    {\parbox[c]{6.6cm}{\includegraphics[width=6.6cm]{images/Campus02_Logo.pdf}} \par}
    \vskip .1em
    {FACHHOCHSCHULE DER WIRTSCHAFT \par}
    \vskip 2em
    {Bachelor Degree Programme \\ Innovation Management \par}
    \vskip 2em
    {by \\ {\Large \textbf{\campusAuthor}} \\ \campusPersKnz \par}
    \vskip 2em
    {supervised and reviewed by \par}
    {\normalsize Dipl.-Ing. Manuela Reinisch \\ FH-Prof. Dipl.-Ing. Dr.techn. Michael Terler \par}
    \vfill
    {
      \begin{flushleft}
        \hfill\parbox[c]{7.3cm}{\includegraphics[width=5cm]{images/Signature}}
        \begin{tabular}{lclcl}
          {\campusCity, \campusDate} & ~\hfill & \hspace*{3cm} & ~\hfill &  \underline{\hspace*{5cm}} \\
          ~\hfill & ~ &   & ~ &  \multicolumn{1}{c}{Signature}   \\
        \end{tabular}
      \end{flushleft}
    }
  \end{center}
  \par
}


\begin{document}

\frontmatter

\renewcommand*{\chapterpagestyle}{scrplain}
\pagestyle{scrplain}

% allow more flexible whitespaces to avoid hyphenation and overfull hboxes
\sloppy

\begin{titlepage}

\subject{\campusSubject}
\title{\campusTitle}

\author{\bfseries \campusAuthor}
\date{\campusCity, \campusDate}

\maketitle

\end{titlepage}
\newpage

\chapter*{Ehrenwörtliche Erklärung}

\begin{otherlanguage}{austrian}
Ich erkläre ehrenwörtlich, dass ich die vorliegende Arbeit selbstständig und ohne fremde Hilfe verfasst, andere als die angegebenen Quellen nicht benützt und die benutzten Quellen wörtlich zitiert sowie inhaltlich entnommene Stellen als solche kenntlich gemacht habe.
\end{otherlanguage}

{
  \begin{flushleft}
    \vskip 2em
    \hfill\parbox[c]{7.3cm}{\includegraphics[width=5cm]{images/Signature}}
    \begin{tabular}{lclcl}
      {\campusCity, \campusDate} & ~\hfill & \hspace*{3cm} & ~\hfill &  \underline{\hspace*{5cm}} \\
      ~\hfill & ~ &   & ~ &  \multicolumn{1}{c}{Signature}   \\
    \end{tabular}
  \end{flushleft}
}

%\chapter*{Acknowledgement}

%Firstly, I would like to express my sincere gratitude to the supernovae for the creation of all substantial atoms and molecules our fantastic world consists of. Their ephemeral power made it possible to form all known and unknown creatures and plants. Without this furtune it would not be possible to live this wunderful life. I thank my family for supporting me spiritually throughout my life and during writing this thesis.

\chapter*{Abstract}

Wireless networking technologies are widely used in communication devices and services in almost every area of daily business. Medical environments, security authorities, and other organizations use them to increase their performance. In particular the Smart Home sector is one of the areas that has been researched extensively. From ecological and technological point of view, the opportunities for new technologies are vast. Therefore, it is important to compare such technologies using ecological parameters. This thesis will give an overview of the wireless network types, ranging from body area to global area networks, and introduces the different Smart Home networks. Furthermore, the major wireless networking technologies of personal area networks, Bluetooth, Bluetooth Smart, Ultra Wideband, ZigBee, Internet Protocol version 6 over Low-Power Wireless Personal Area Network and Wi-Fi, are discussed in detail. A comparison is made with current chipset manufacturers implementations to create a tabular overview of these technologies and their suitability in Smart Home networks. Results show that every technology has its optimal field of application in a modern Smart Home. However, it can be concluded that further experiments will show more accurate results.

\vspace{2cm}

\textbf{Keywords:} \campusKeywords

\chapter*{Zusammenfassung}

\begin{otherlanguage}{austrian}
Drahtlose Netzwerktechnologien zählen zum Rückgrad der Kommunikationseinrichtungen unseres täglichen Arbeitsalltags. In medizinischen Einrichtungen, Sicherheitsbehördern und anderen Organisationen werden diese Technologien zur Leistungssteigerung und zur Erhöhung der Effizienz eingesetzt. Zunehmend wird in letzter Zeit auch der Smart Home Sektor Ziel intensiver Forschung im Bereich der Netzwerktechnologien. Aus ökologischer und technologischer Sicht ergeben sich Chancen für neue effizientere Technologien. Daher ist es wichtig, solche Technologien unter den Gesichtspunkten ökologischer Parameter zu vergleichen. Diese Arbeit gibt einen Überblick über die drahtlosen Netzwerktypen, die von körpernahen bis hin zu globalen Netzwerken reichen und erläutert die verschiedenen Smart Home Netzwerke. Aus dem Bereich der persönlichen Netzwerke werden die wichtigsten drahtlosen Netzwerktechnologien vorgestellt. Dabei werden Bluetooth, Bluetooth Smart, Ultra-Breitband, ZigBee, Internet Protocol Version 6 für Wireless Personal Area Network mit niedrigem Energieverbrauch und Wi-Fi ausführlich diskutiert. Ein Vergleich mit Implementierungen von Chip-Herstellern führt zu einer tabellarischen Übersicht dieser Technologien und schafft die Grundlage für die Beurteilung der Eignung in Smart Home Netzwerken. Die Ergebnisse zeigen, dass jede Technologie ihr optimales Einsatzgebiet in einem modernen Smart Home hat. Es kann jedoch der Schluss gezogen werden, dass künftige Experimente zur Eignung dieser Technologien im Smart Home diese Beurteilung vervollständigen können.
\end{otherlanguage}

\vspace{2cm}

\textbf{Keywords:} \campusKeywords

\mainmatter

\pagenumbering{roman}
\setcounter{page}{1}

\pagestyle{scrheadings}
\setlength{\parindent}{0em}
\renewcommand*{\chapterpagestyle}{scrheadings}

\tableofcontents

\clearpage

\pagenumbering{arabic}
\setcounter{page}{1}

\chapter{Introduction}
\label{ch:intro}

This \LaTeX{} template is designed for the creation of thesis documents (bachelor, master, phd) and targets both beginner and experienced users of \LaTeX{}. It supports all basic functionality and requirements of a technical document such as the inclusion of graphics, math, tables, references, bibliography and much more. In contrast to a standard LaTeX document this template not only loads all state of the art packages (\path{scientific.cls}) to provide the best functions for each task, but also includes a separate document for the style/layout of the document (\path{style.tex}). It therefore tries to separate functionallity and layout as much as possible. And the best, everything is documented in the code.

\chapter{Wireless Networking Technologies}
\label{ch:wireless-networking}

Everyday we use networking technologies to exchange information. The \gls{ict} industry is the innovation and development engine of these technologies. They allow us to use information technology wherever and whenever we need them. Whether it is in the office, at home, or on the way, \gls{ict} is available everywhere. \footcite[Cf.][1]{Pahlavan2009}

There are many different technologies used to keep us connected to the global Internet by wire and wireless. More than 300 submarine cables are installed globally to interconnect all regions of the world using high-speed and high-capacity data routes. A vast number of access points builds the gap between the wired and the wireless \gls{ict}. The technologies used for wireless networks are commonly differentiated by their geographic range of operation. \footcite[Cf.][49--52]{Ephremides2013}

\begin{figure}[ht]
  \centering
  \includegraphics[
    width=\textwidth,
  ]{images/wireless_network_types}
  \caption{Overview of the Wireless Network Types, adapted from: \cite[6]{Latre2011}}
  \label{fig:wireless-network-type}
\end{figure}

\Cref{fig:wireless-network-type} shows an overview of the wireless network types, sorted in ascending order of their communication range. They can be categorized into \gls{wban}, \gls{wpan}, \gls{wlan}, \gls{wcan}, \gls{wman}, \gls{wwan}, \gls{wran}, and \gls{wgan}. \footcite[Cf.][6]{Latre2011}

\section{Wireless Body Area Network}

Regarding their geographical coverage, the smallest is the \gls{wban}. A \gls{wban} covers an area of about one meter and has a very limited amount of throughput, which is the domain of health monitoring applications. Devices operating in a \gls{wban} are for example wearables, attached to the body. \footcite[Cf.][1660]{Movassaghi2014}

The IEEE 802.15.6 standard define Physical and Medium Access Control layers for \gls{wban}. \footcite[Cf.][1]{Kwak2010}

\begin{figure}[ht]
  \centering
  \includegraphics[
    height=4cm,
    width=7cm,
  ]{images/wban}
  \caption{Typical \gls{wban} scenario, source: Own diagram}
  \label{fig:wban}
\end{figure}

Health and fitness applications use biosensors measuring biochemical interactions with the human body. Those signals are transmitted and monitored remotely. \gls{wban} sensors, as shown in \cref{fig:wban}, allow users to operate freely without wires. This gives the individual maximum amount of mobility. A typical \gls{wban} sensor system consists of a biosensor, a transmission unit, and a control and processing unit. They are used to monitor fitness or disease progression using typical parameters: \footcite[Cf.][69--70]{Minoli2013}

\begin{itemize}
  \item Temperature: Both the user and ambient temperature are measured by specific sensors and recorded for further processing and evaluation.
  \item Blood Glucose Level: This is an important parameter for e.g. diabetes patients.
  \item Blood Pressure: A simple indicator for a person's health.
  \item Oxygen Saturation: Measured indirectly, which means it cannot measure the amount of oxygen used by the user.
  \item Heart Rate: Increasingly important for top athletes.
  \item Step Count: An important parameter to measure daily fitness.
\end{itemize}

Standardization of \gls{wban} technologies is an important task of several research groups. Nowadays a number of technologies are used to operate within a \gls{wban}: \footcite[Cf.][71]{Minoli2013}

\begin{itemize}
  \item ZigBee, an IEEE standard based protocol, enabling the use of low power wireless technology with less implementation effort.
  \item Bluetooth Smart, is the low power variant of the Bluetooth protocol, also an IEEE standard. The data rate is limited, but battery life make it suitable for medical applications in \gls{wban} systems.
  \item NFC, a contactless communication technology used mainly by smart phones or tablets to gather information of compatible sensors.
\end{itemize}

Bluetooth and ZigBee are mainly used in \gls{wpan} environments and adopted for applications operating using a \gls{wban}. The IEEE 802.15.6 task group was founded to standardize \gls{wban} technologies. \footcite[Cf.][10]{Latre2011} The current standard proposes different frequency layers for data transmission: The \gls{nb} ranging from 400 MHz to 2.4 GHz, the \gls{uwb} operating between 3.1 GHz and 11.2 GHz and the \gls{hbc} which uses 10 MHz to 50 MHz. Another technology is recently been added to the \gls{wban} technology list: low-power WiFi, which is an adaptation of the IEEE 802.11 standard. \footcite[Cf.][]{Kwak2010}\footcite[Cf.][180--187]{Minoli2013}

\section{Wireless Personal Area Network}

The \gls{wpan} has an indoor coverage up to 20 meters and is capable of high data rates up to 110 Mbps. \footcite[Cf.][676]{Garg2007} The IEEE 802.11ah low-power WiFi protocol is currently under standardization, but other protocols like \gls{rfid}, ZigBee, Bluetooth or \gls{nfc} have already proved their concepts and are widely used in \gls{wpan} environments. They connect \gls{wban} devices to higher level network architectures like \gls{wpan} or \gls{wlan}. \footcite[Cf.][]{wban_wpan_wlan}

\begin{figure}[ht]
  \centering
  \includegraphics[
    width=7.5cm,
  ]{images/wban_wpan_wlan}
  \caption{Interconnection of \gls{wpan}, source: \cite{wban_wpan_wlan}}
  \label{fig:wban_wpan_wlan}
\end{figure}

\Cref{fig:wban_wpan_wlan} shows an overview of a typical \gls{wpan} architecture. The \gls{wpan} is interconnecting devices and services from the \gls{wban}. Body sensors and wearables of the \gls{wban} are enabled to transmit their gathered information during their monitoring process to services or devices in \gls{wpan} for storage or further processing. The \gls{wpan} is focused on interconnecting personal equipment for portable and mobile communication. \footcite[Cf.][654-704]{Garg2007}

Such devices operating in a \gls{wpan} are: \footcite[Cf.][654-704]{Garg2007}

\begin{itemize}
  \item \glspl{pc} or Notebooks: Desktop \glspl{pc} and notebooks are used for web surfing, home office, gaming and other applications. In conjunction with the \gls{wban} devices they have mainly the role of a storage and control device, and they are able to bridge between various wired or wireless networks.
  \item \glspl{pda} or Tablets: Nowadays a very important device to control entertainment systems, act within social networks and other mobile applications.
  \item Mobile Phones: Maybe the only technical device spread across the whole world and de facto the standard device which every human being is using every single day. Most mobile phones have multiple wireless technologies implemented, a wireless cellular network technology and a \gls{wban} / \gls{wpan} technology, like Bluetooth.
  \item Printers: Many printers, both standalone and mobile models, are capable of accepting print jobs through \gls{wban} technologies, like Bluetooth and NFC.
  \item Speakers and Microphones: Bluetooth audio devices, such as hands-free equipment, are widely used for fitness and entertainment.
  \item Consumer Electronics: Many remote controls are based on Bluetooth or other \gls{wpan} technologies to control electronic equipment and for streaming audio and video signals.
\end{itemize}

\begin{figure}[ht]
  \centering
  \includegraphics[
    width=13cm,
  ]{images/wpan}
  \caption{Overview of \gls{wpan} standards, source: \cite[143]{Jawad2014}}
  \label{fig:wpan}
\end{figure}

Nowadays the social networks are used to share arbitrary personal information of different media types around the world. These media consists of pictures taken by mobile phones, videos shot by action cams and many others. Those media streams have different requirements to \gls{wpan} technologies, such as bandwidth, data throughput and robustness. \footcite[Cf.][60]{Yang2008}

\gls{uwb} and IEEE 802.11ah are exposed to provide high date rates up to 1 Gbps for streaming high definition video signals. Conversely, many other applications demand on lower data rates for transmission of monitoring and control commands, the domain where ZigBee is located. \Cref{fig:wpan} gives a summary of the IEEE \gls{wpan} technology standards. \footcite[Cf.][654]{Garg2007}

\section{Wireless Local Area Network}

The IEEE 802.11 working group is specialized on defining standards for the \gls{wlan}. The data rate is extended to 700 Mbps and the indoor coverage is around 50 meters, whereas the outdoor range is up to approximately 250 meters. The IEEE 802.11ad and IEEE 802.11ay groups are working on standardization of very fast 60 GHz communication offering data rates beyond 5 Gbps for IEEE 802.11ad respectively 100 Gbps for IEEE 802.11ay. Key requirements for \gls{wlan} technologies are: \footcite[Cf.][]{Bellalta2015}

\begin{itemize}
  \item Coexistence: \glspl{wlan} operate within restricted \gls{ism} bands and therefore the \gls{wlan} technologies have to be able to coexist with other wireless networks operating at their bands.
  \item Throughput: To further increase the throughput of a \gls{wlan} system, the use of multiple antennas and an optimized or new wireless technology has to be discovered.
  \item Energy Efficiency: This is a very difficult requirement which is in contrast to throughput increase. Nonetheless is required to further decrease the energy consumption of the \gls{wlan} systems.
  \item Backward Compatibility: Although it leads to inefficiency, a mechanism to provide backward compatibility to former IEEE 802.11 standards should be provided.
\end{itemize}

\begin{figure}[ht]
  \centering
  \includegraphics[
  width=13.5cm,
  ]{images/wlan}
  \caption{Typical \gls{wlan} scenario, source: \cite{devolo}}
  \label{fig:wlan}
\end{figure}

\Cref{fig:wlan} shows a common indoor \gls{wlan} environment spread across multiple rooms. A smart TV streaming high definition video signals and media meta information from the Internet, a tablet receiving the latest news and social media updates, and a notebook running a fast-paced online game. The requirements to \gls{wlan} technologies are very diverse and span from low rate secure data transmission to very high rate data streaming. \footcite[Cf.][]{devolo}

\begin{table}[h!]
  \renewcommand{\arraystretch}{2}
  \newcolumntype{P}[1]{>{\raggedright}p{#1}}
  \centering
  \begin{tabular}{ P{3cm} P{2.8cm} P{2.8cm} P{2.8cm} >{\arraybackslash}P{2.8cm} }
    \toprule
    \rowcolor{Gainsboro}
    \textbf{Feature / IEEE standard} & \textbf{802.11b} & \textbf{802.11g/a} & \textbf{802.11n} & \textbf{802.11ac} \\
    \midrule
    \textbf{Maximum data rate per stream} & 11~Mbps & 54~Mbps & >100~Mbps & >500~Mbps / >1000~Mbps \\
    \textbf{Frequency band} & 2.4~GHz & 2.4~GHz / 5~GHz & 2.4~GHz and 5~GHz & 5~GHz \\
    \textbf{Channel width} & 20~MHz & 20~MHz / 20~MHz & 20~MHz and 40~MHz & 20~MHz, 40~MHz, 80~MHz, 160~MHz and 80+80~MHz \\
    \textbf{Antenna technology} & \gls{siso} & \gls{siso} & \gls{mimo} & \gls{mimo} / \gls{mumimo} \\
    \textbf{Transmission technique} & \gls{dsss} & \gls{dsss} and \gls{ofdm} & \gls{ofdm} & \gls{ofdm} \\
    \textbf{Maximum number of spatial streams} & 1 & 1 & 4 & 8 \\
    \textbf{Beamforming-capable} & No & No & Yes & Yes \\
    \bottomrule
  \end{tabular}
  \caption{Overview of IEEE 802.11 standards, adapted from: \cite[85]{Bejarano2013}}
  \label{table:wlan}
  \vspace{-1.5em}
\end{table}

\Cref{table:wlan} shows a comparison of the IEEE 802.11 standards. With IEEE 802.11ac the data rate massively increases. This is due to the modulation which is used by this standard and the ability to enable simultaneous multiple user connections. \footcite[Cf.][84--85]{Bejarano2013}

\section{Wireless Campus Area Network}

A \gls{wcan} is composed by connecting multiple \gls{wlan} across offices or throughout entire campus-like buildings. Because of the geographical area used by campuses, the \gls{wcan} sometimes uses \gls{wman} technologies to provide good wireless coverage over the campus area. The IEEE 802.11 and IEEE 802.16 standards are used for high data rate wireless networks with good indoor coverage. Using \gls{mimo} and \gls{mumimo} technology, the number of simultaneous connections of a wireless access point can be increased. \footcite[Cf.][]{Chen2009}

\section{Wireless Metropolitan Area Network}

The increasing demand for high data rate wireless transmissions in urban and sub-urban areas, requires new technologies. Covered by the IEEE 802.16 family of standards, the \gls{wman} has a geographical coverage of operation of about 10 km. IEEE 802.16 based technologies are used to operate large-scale networks across cities.\footcite[Cf.][713]{Garg2007} The IEEE 802.16 standard is called \gls{wimax}.\footcite[Cf.][1]{Kirmse2009} The \gls{wimax} protocol family allows up to 64,000 simultaneous connections per access point. Two types of \gls{bwa} services are currently available: \footcite[Cf.][195--249]{Ephremides2013}

\begin{itemize}
  \item Fixed Broadband: This type attempts to provide the same services as traditional wired broadband connections (e.g. \gls{dsl} or cable modems).
  \item Mobile Broadband: Nomadicity and mobility are the features of this type of \gls{bwa} services type. Nomadicity is the ability to connect to the \gls{bwa} base station from different locations within the coverage range of a \gls{bwa} base station. Mobility describes the feature of keeping connections to the \gls{wman} active while moving.
\end{itemize}

In 2011 the \gls{wimax} Forum had stated a number of 583 \gls{wimax} installations in over 150 countries worldwide. \footcite[Cf.][246]{Ephremides2013}

\section{Wireless Wide Area Network}

The standards defining technologies for \gls{wwan} are typically based on cellular network topologies. A \gls{wwan} covers a geographical range of up to 10km. Technologies in \gls{wwan} are typical 2G and 3G systems: \footcite[Cf.][402--403]{Ephremides2013}

\begin{itemize}
  \item \gls{gsm} (2G): The \gls{gsm} standard is based on \gls{tdma} and can achieve data rates up to 1 Mbps using \gls{edge} technology.
  \item \gls{umts} (3G): This standard is based in the \gls{cdma} and provides up to 42 Mbps data transfer rates when \gls{hspa+} can be used.
\end{itemize}

The IEEE 802.16 standard defines a mobile broadband technology called Mobile \gls{wimax} which is a fully compliant 3G system at data rates up to 80 Mbps.\footcite[Cf.][2-3]{Kirmse2009} The main difference to the fixed type \gls{wimax} is the support for mobile devices, such as mobile phones, tablets or laptop computers. Embedded Mobile \gls{wimax} solutions or Mobile \gls{wimax} extension cards have to be used in conjunction with omni-directional antennas. \footcite[Cf.][3--4]{Kirmse2009}

\section{Wireless Regional Area Network}

The IEEE 802.22 working group defines standards for cognitive radio-based PHY/MAC/air interfaces for the use in \gls{wran}, which have a typical geographical coverage of 100km. Their target is to provide \gls{bwa} with data rates up to 1.5 Mbps to areas with sparse Internet access. The spectrum of IEEE 802.22 is located in the TV Broadcasting 54 MHz to 862 MHz band. \footcite[Cf.][2]{Romme2008}

\section{Wireless Global Area Network}

\begin{figure}[ht]
  \centering
  \includegraphics[
  width=10cm,
  ]{images/all_ip_network}
  \caption{Overview of the all-ip-network architecture of Smart Grids, source: \cite{mulligan}}
  \label{fig:all_ip}
\end{figure}

The \gls{wgan} consists of satellite-based transmission systems spanning around the world at the near-earth orbit. Additionally, the upcoming 5G wireless network standards will provide wireless communication globally. It tends to create a \gls{wwww}. Based on technologies like \gls{cdma}, \gls{ofdm}, \gls{uwb} and IPv6, it will provide All-IP wireless network systems. \Cref{fig:all_ip} gives an overview of the architecture of future All-IP wireless networking. The network is clustered into smaller Smart Grids which provide local network access and allow a hand-off of active connections through Smart Grids. \footcite[Cf.][]{Lin2005}

\chapter{Smart Home}
\label{ch:smart-home}

Smart Homes are interactive spaces where consumer electronics, healthcare products, security installations, energy monitoring, mobile phones, home computing equipment and many more electronic devices and services are interconnected to provide remote control and monitoring services. \footcite[Cf.][]{Alam2012}

The Smart Home is based on pervasive or ubiquitous computing, a model where the user interacts in a normal manner with his home. This computing model follows the trends of evolution in information technology. Developing from communication devices, over mobility, past ubiquitous systems, hence to ambient intelligence, the current Smart Home systems are artificial computing devices. \footcite[Cf.][]{Lalanda2010}

Those devices have to follow the aspects of miniaturing, to ideally disappear at their operating location. Furthermore they have to be able to communicate with a whole bunch of other Smart Home systems. The need to be autonomous in the context of energy supply. Batteries and energy harvesting are key technologies to their energy autonomy. \footcite[Cf.][]{Lalanda2010}

Technologies from different network types are found in or incorporating with Smart Home systems. Current Smart Home systems are communicating through wireless networking technologies. Ranging from very short range \gls{wban} to wide range \gls{wwan}. They form according to their connected devices and services different types of networks. \footcite[Cf.][]{Saito2013}

The following chapters discusses the main types of Smart Home networks.

\section{Wireless Home Entertainment Networks}

The \gls{when} is the domain of streaming rich media content from centralized services or mobile devices. Nowadays video streaming is a booming market. Amazon Video, Netflix, Youtube and DailyMotion are some examples of streaming video hosters. They offer full high definition videos, movies and clips. Several million users are watching streaming content every day. Such media streams are often transported over the \gls{http} using a container for the video and audio channels. There are several embedded processor manufacturer providing processors with hardware-based acceleration for decoding several video and audio formats. \footcite[Cf.][593--596]{Kurose2012}

The NVIDIA Tegra X1 mobile processor is a powerhorse featuring 4K (4096 x 2160 pixels) video codec. With an operating power of less than 10 Watts this processor is designed for video intensive tasks while operating fanless on embedded platforms. \footcite[Cf.][]{tegrax12015}

Such embedded platforms are data sinks for streaming media content. The data sources are commonly located on centralized storage devices or streamed directly from the Internet or mobile devices. Embedded platforms featuring high performance processors are also data sources for media streams. \footcite[Cf.][123]{Augusto2006}

The wireless networking technologies mapping home entertainment networks must provide sufficient bandwidth for streaming of high definition media. A 4K movie has a bitrate of around 40 Mbps. Furthermore the wireless network has to provide low latency and high reliability. In case of a packet loss, the missing packets needs to be retransmitted within a few milliseconds. Additionally, multipath and coexistence issues needs to get addressed by \gls{when} technologies. \footcite[Cf.][]{arris2014}

\section{Wireless Home Automation Networks}

Wireless sensors and actuators are the devices forming \glspl{whan}. Remote applications are used to monitor and control those devices. \gls{whan} devices use low-power wireless networking technologies and they can operate energy autonomously if they are battery-powered or use energy harvesting mechanisms. The typical use cases for \gls{whan} are spread across all areas of a users home: \footcite[Cf.][]{Gomez2010}

\begin{itemize}
  \item Light control: Through presence detection or luminous intensity measuring sensors.
  \item Remote Control: By using intelligent \gls{rf} sensors.
  \item Smart Energy: A combination of several sensors could detect the energy needs and trigger control systems.
  \item Remote care: Wearables can monitor health parameters from patients and transmit them to receiver nodes.
  \item Security and safety: Any sensor might be a source for security application systems.
\end{itemize}

According to the use cases for \glspl{whan} the number of sensors and actuators is resulting in a very dense network. The multipath characteristics of a home and the interference with other wireless networks are issues to \glspl{whan}. Additionally the nodes within a \gls{whan} should be able to retransmit control messages to other nodes, building a multi-hop end-to-end communication facility. \footcite[Cf.][]{Gomez2010}

\section{Wireless Machine-To-Machine Networks}

The number of \gls{m2m} networks is growing exponentially and building the \gls{iot}. Sensors and other intelligent devices are connected to the Internet through home gateways. \gls{m2m} is very close to \gls{whan} in case of the network nodes. Both networks consists of sensors, actuators and remote control and monitoring devices and applications. The main difference between \gls{m2m} and \gls{whan} is that \gls{whan} consists of \gls{m2m}, \gls{h2m} and \gls{mih} nodes. \footcite[Cf.][]{Jara2013}

The typical network characteristics and requirements of \gls{m2m} devices are: \footcite[Cf.][]{Minoli2013}

\begin{itemize}
  \item Retransmission: Recover from packet loss and other transmission issues.
  \item Scalability: Support several hundreds of nodes at various link speeds.
  \item Prioritization: Routing of emergency messages at higher priority.
  \item Security: Prevent remote attacks through strong encryption algorithms.
\end{itemize}

\section{Wireless Home Networks}

Embedded processors can be found in our mobile phones, laptops, watches, light bulbs, and many more. The more intelligent those devices get, the more information they will provide. Therefore a ubiquitous networking environment is essential to transmit this information. The previous network types can be summarized to the \glspl{whn}, whereas sensors, actuators, consumer electronics, smart phones, \glspl{wpan} or \glspl{wlan}, and even the electric car are part of. \Cref{fig:home_network_parameters} summarizes the characteristics and requirements of the different network types and their nodes. \footcite[Cf.][7286--7290]{Mendes2015}

\begin{figure}[ht]
  \centering
  \includegraphics[
    width=8cm,
  ]{images/han_parameters}
  \caption{Parameters of a \gls{han}, adapted from: \cite[7290]{Mendes2015}}
  \label{fig:home_network_parameters}
\end{figure}

\chapter{Results}
\label{ch:results}

\begin{figure}[ht]
	\centering
	\includegraphics[width=5cm]{images/Placeholder}
	\caption{Typical figure, source: Own diagram}
	\label{fig:1}
\end{figure}

\chapter{Key Parameters}
\label{ch:key-params}

In ubiquitous computing of Smart Home networks, the different wireless networking protocols do not perform well in any Smart Home scenario. Therefore, the wireless networking technology which performs best in the specific Smart Home application should be choosen. To perform a comparison of wireless networking technologies in accordance to different Smart Home applications, a wireless protocol should address the following issues: \footcite[Cf.][7289--7293]{Mendes2015}

\begin{itemize}
  \item Energy efficiency: The low power consumption of the transmission layers.
  \item Bandwidth efficiency: Spread spectrum efficiency, channel width, and coexistence stability.
  \item Scalability: The number of nodes supported in a network.
  \item Transmission Range: The theoretical coverage range and the penetration of obstacles.
  \item Throughput: The maximum rate of transmission.
  \item Latency: Theoretical latency between two nodes.
  \item Mobility: Efficient mobile node management.
\end{itemize}

The taxonomy of parameters for the wireless networking technologies is shown in \cref{fig:taxonomy_of_parameters}. The evaluation of the parameters is discussed in \cref{ch:eval}.

\begin{figure}[ht]
  \centering
  \includegraphics[
    width=12cm,
  ]{images/taxonomy_of_parameters}
  \caption{Taxonomy of parameters for evaluation of wireless networking technologies, adapted from: \cite[7290]{Mendes2015}}
  \label{fig:taxonomy_of_parameters}
\end{figure}

\chapter{Evaluation and Comparison}
\label{ch:eval}

\begin{table}[h!]
  \renewcommand{\arraystretch}{2}
  \newcolumntype{P}[1]{>{\raggedright}p{#1}}
  \centering
  \begin{tabular}{ P{2.3cm} P{1.8cm} P{1.8cm} P{1.8cm} P{2cm} P{2cm} >{\arraybackslash}P{1.8cm} }
    \toprule
    \rowcolor{Gainsboro}
    \textbf{Parameters} & \textbf{Bluetooth} & \textbf{Bluetooth Smart} & \textbf{UWB} & \textbf{ZigBee} & \textbf{6LoWPAN} & \textbf{Wi-Fi} 802.11ad \\
    \midrule
    \textbf{Chipset} & CC2560 \footnotemark & CC2640 \footnotemark & {PL3120 \footnotemark \\ PL3130 \footnotemark} & CC2630 \footnotemark & CC2630 \footnotemark & Wil6200 \footnotemark \\
    \textbf{Frequency Band} & 2.4~GHz to 2.48~GHz & 2.4~GHz to 2.48~GHz & 3.1~GHz to 5.8~GHz & 868.3~MHz / 902~MHz to 928~MHz / 2.4~GHz to 2.48~GHz & 868.3~MHz / 902~MHz to 928~MHz / 2.4~GHz to 2.48~GHz & 57.24~GHz to 65.88~GHz \\
    \textbf{Number of Channels} & 79 & 40 & 3 & 27 & 27 & 4 \\
    \textbf{Channel Bandwidth} & 1~MHz & 2~MHz & 528~MHz & 1~MHz / 2~MHz / 5~MHz & 1~MHz / 2~MHz / 5~MHz & 2.16~GHz \\
    \textbf{Modulation} & \gls{gfsk} & \gls{gfsk} & \gls{ofdm} & \gls{dsss} & \gls{dsss} & \gls{ofdm} \\
    \textbf{Bandwidth Spreading} & \gls{fhss} & \gls{fhss} & \gls{dsss} & \gls{dsss} & \gls{dsss} & \gls{ofdm} \\
    \textbf{Max. Data Rate} & 2.1~Mbps & 0.27~Mbps & 675~Mbps & 20~kbpss / 40~kbps / 250~kbps & 20~kbps / 40~kbps / 250~kbps & 6.75~Gbps \\
    \textbf{Transmission Range} & 10~m & 100~m & 10~m & 10~m & 10~m & 10~m \\
    \textbf{Energy Consumption} & 1~W & 500~mW & 330~mW & 20~mW & 20~mW & 700~mW \\
    \textbf{Network Type} & Piconet / Scatternet & Scatternet & Piconet & Scatternet & Piconet & Peer-to-Peer \\
    \textbf{Max. Number of Nodes} & 7 & 5917 & 16 & 1024 & 100 & 1 \\
    \textbf{Latency} & 100~ms& 6~ms & <1~ms & 100~ms & 100~ms & 10~{\textmu s} \\
    \textbf{Mobility} & Yes & Yes & Yes & Yes & Yes & No \\
    \bottomrule
  \end{tabular}
  \caption{Parameters of \gls{wpan} technologies, source: Own diagram}
  \label{table:parameters}
  \vspace{-1.5em}
\end{table}
\footcitetext[Cf.]{cc2560}
\footcitetext[Cf.]{cc2640}
\footcitetext[Cf.]{pl3120}
\footcitetext[Cf.]{pl3130}
\footcitetext[Cf.]{cc2630}
\footcitetext[Cf.]{cc2630}
\footcitetext[Cf.]{wil6200}

The parameters of the wireless networking technologies, shown in \cref{table:parameters}, were researched from chipset manufacturers and combined with the information of \cref{ch:wpan}. Furthermore, the following comparison can be conducted as illustrated in \cref{table:comparison}.

The applicability of the Bluetooth technology is far away from the expectations. Because of the low data rate, Bluetooth is not usable for \gls{when} applications. The reason for the ordinary rating in the other \gls{whn} categories, is due to the small number of network nodes.

A totally different picture shows Bluetooth Smart. Caused to the low data rate, the Bluetooth Smart technology is not well suitable for \gls{when} applications. But Bluetooth Smart is a good choice for \gls{whan} and wireless \gls{m2m} networks. Both network types have the requirements of very low power consumption, a scalable network architecture and the support for hundreds a network nodes. Bluetooth Smart offer a good coverage range and is also applicable for outdoor sensors. But Bluetooth Smart is not very easily integratable into heterogenous \gls{whn} environments, due to the non-\gls{ip} stack.

\begin{table}[h!]
  \renewcommand{\arraystretch}{2}
  \newcolumntype{P}[1]{>{\raggedright}p{#1}}
  \centering
  \begin{tabular}{ P{3cm} P{2.5cm} P{2.5cm} P{2.5cm} >{\arraybackslash}P{2.5cm} }
    \toprule
    \rowcolor{Gainsboro}
    \textbf{Wireless Networking Technologies} & \textbf{Wireless Home Entertainment Network} & \textbf{Wireless Home Automation Network} & \textbf{Wireless \gls{m2m} Network} & \textbf{Wireless Home Network} \\
    \midrule
    \textbf{Bluetooth} & Low & Medium & Medium & Medium \\
    \textbf{Bluetooth Smart} & Low & High & High & Medium \\
    \textbf{UWB} & High & Medium & Medium & Low \\
    \textbf{ZigBee} & Low & High & High & Medium \\
    \textbf{6LoWPAN} & Low & High & High & High \\
    \textbf{Wi-Fi 802.11ad} & High & Low & Low & Medium \\
    \bottomrule
  \end{tabular}
  \caption{Aptitude of \gls{wpan} technologies in Smart Home networks, source: Own diagram}
  \label{table:comparison}
  \vspace{-1.5em}
\end{table}

The \gls{uwb} technology is very well suited for \gls{when}. \gls{uwb} offer a high data rate of several hundreds of Mbps. The technology is very robust to interference and has low power consumption. \gls{uwb} is a versatile technology, just the number of network nodes is not large enough to be the first choice in \gls{whan} and wireless \gls{m2m} networks. The non-\gls{ip} stack is also a barrier for heterogenous \gls{whn} environments.

ZigBee is the technology for the \gls{iot}, therefore the rating is high in the respective categories, \gls{whan} and wireless \gls{m2m} networks. Low power consumption and a large number of network nodes are the main requirements for those network types. ZigBee does not offer high data rates, nor is the protocol not \gls{ip} conform. Therefore, ZigBee is not a good choice for \gls{when} or heterogenous \gls{whn} environments.

The next candidate makes that possible, \gls{6lowpan}: Easy integratable into \gls{ip} networks, low power consumption and a high maximum number of network nodes. \gls{6lowpan} is well suited for \gls{whan}, wireless \gls{m2m} networks and heterogenous \gls{whn} environments. Because of the low data rate, \gls{6lowpan} is not very attractive to \gls{when}.

The escapee from the Wi-Fi technologies, the IEEE 802.11ad standard, is very well suited for high speed data transmissions as needed in \glspl{when}. But, the number of simultaneous data transfers to different network nodes is very limited. That is why the IEEE 802.11ad is not applicable in \gls{whan} or wireless \gls{m2m} networks. Due to the fact that IEEE 802.11ad \gls{mac} layer is combinable with other IEEE 802.11 standards in multi-\gls{mac} chipsets, the IEEE 802.11ad is integratable in heterogenous \gls{whn} environments.

\chapter{Conclusion and Future Work}
\label{ch:conclusion}

Due to the wide spread of wireless networking technologies along most domains of our life, more and more attention has been focused on ecological parameters of these technologies. The main challenge from the many difficulties involved in the selection of wireless networking technologies is in providing a guidance to the aptitude for a specific domain. It is necessary for a large number of wireless applications to provide an energy efficient overall picture, because power consumption is a crucial parameter for decisions on selecting wireless products and services.

The sector of ecologically designed Smart Homes is part of the concentrated research on energy efficient wireless networking technologies. Because of the many and various wireless networking types in Smart Homes, the requirements for applications and services in Smart Homes spread over multiple parameters.

This thesis presents a guidance for the selection of wireless networking technologies in Smart Homes based on their standards and implementations by various manufacturers. The first part of this thesis introduces the different wireless network types according to their coverage range and gives an overview to the wireless Smart Home networks. Examples are given, taken from their respective area of application and the key issues of wireless Smart Home networks are introduced.

The subject of the second part of this thesis is the introduction of probably suitable wireless networking technologies for energy efficient Smart Home applications. In order to answer the research question I reviewed six energy efficient wireless networking technologies used in \glspl{pan}. Those technologies are all open standards, and were developed and maintained by different alliances.

The discussed \gls{wpan} technologies can be found in Smart Home applications in every corner and their possibilities are vast. All of these technologies are promising within their domain but Bluetooth will be less significant. Especially the network topology is a downside of Bluetooth. \gls{uwb} and IEEE 802.11ad offer extreme high data rates, and ZigBee and \gls{6lowpan} are both technologies for larger sensor and actuator networks.

The comparison of the \gls{wpan} technologies in conjunction with their aptitude for different Smart Home applications shows that every technology has its hotspots and right to exist in the Smart Home networking environment.

It would be interesting to investigate the aptitude of the \gls{wpan} technologies with experiments in simulated environments. Such simulations could deliver accurate results, which completes this guidance by practical experience with the focus on energy-efficiency, security and coexistence. These parameters are meaningful areas for further study, to adapt and develop new Smart Home products.


% uncomment to include all literature sources
%\nocite{*}

\clearpage

\ihead[]{References}
\printbibheading[heading=bibintoc, title=References]
\begin{bibhyphenation}
  \printbibliography[heading=subbibintoc, title=Bibliography, filter={printed}]
  \printbibliography[heading=subbibintoc, title=Webography, filter={online}]
\end{bibhyphenation}

\clearpage

\ihead[]{\leftmark}
\listoffigures
\listoftables
\printglossary[type=\acronymtype]

\end{document}